% Book 4, Year 2, chapter 7 -- Enums, pattern matching and error handling.
% Provenance (never printed): sources/book4/BRIEFS.md, chapter 7.
\chapter{Enums, Pattern Matching and Error Handling}\label{ch:b2:rust-enums-and-errors}

The weather stations write their temperatures with one decimal. A C function reads a field into tenths of a degree,
and returns $-1$ when the field is not a number, the usual C convention.

\omcode{code/bachelor-2/07-rust-enums-and-errors/listings/temp_code.c}{5}{12}{An error code inside the range of the
answers}\label{lst:b2:rust-enums-and-errors:temp-code}

\omoutput{code/bachelor-2/07-rust-enums-and-errors/outputs/temp_code.txt}{$-0.1$\,\textdegree C and an empty field
give the same answer.}\label{out:b2:rust-enums-and-errors:temp-code}

A temperature of $-0.1$\,\textdegree C is indistinguishable from a missing one. The function has two kinds of
answers, a temperature or a failure, and its return type, \code{int}, has room for only one kind: the failure
borrowed a value from the temperatures. This chapter gives each kind of answer its own place in the type, and makes
the compiler check that every caller looks at both.

\begin{recall}
Book 3: a C \code{union} stores one of several types in the same bytes, and nothing records which; a C \code{enum}
names integer constants. An inductive set is defined by base elements and constructors, and a property of all its
elements is proved by structural induction: for the base elements, then for each constructor assuming it for the
parts. \cref{ch:b2:rust-ownership}: \omterm{def:b2:rust-ownership:ownership}{ownership}, \code{Box}, and the \omterm{def:b2:rust-enums-and-errors:panic}{panic} that stops a program on a bug.
\end{recall}

\section{Product and sum types}\label{sec:b2:rust-enums-and-errors:types}

\begin{definition}[Product type, sum type, algebraic data type]\label{def:b2:rust-enums-and-errors:adt}
A \emph{product type}\index{product type} holds one value of each of its component types at once: a tuple
\code{(A, B)} or a \code{struct}. A \emph{sum type}\index{sum type} holds one value of exactly one of its component
types, together with which one: a Rust \code{enum}. The types built from basic types by products and sums, possibly
recursively, are the \emph{algebraic data types}\index{algebraic data type}.
\end{definition}

\begin{definition}[Enumeration, enum variant]\label{def:b2:rust-enums-and-errors:enum}
A Rust \emph{enumeration}\index{enumeration (Rust)} \code{enum E { V1(...), V2 { ... }, V3 }} lists its \emph{enum
variants}\index{enum variant}, each a name with its own fields, unnamed, named or none. A value of type \code{E} is
built by one variant, \code{E::V1(...)}, and records which.
\end{definition}

\omcode{code/bachelor-2/07-rust-enums-and-errors/listings/shapes_enum.rs}{4}{31}{Three shapes, one
type}\label{lst:b2:rust-enums-and-errors:shapes}

\omoutput{code/bachelor-2/07-rust-enums-and-errors/outputs/shapes_enum.txt}{The derived \code{Debug} prints each
variant with its fields.}\label{out:b2:rust-enums-and-errors:shapes}

The names \emph{product} and \emph{sum} come from counting.

\begin{proposition}[Counting the values of algebraic types]\label{prop:b2:rust-enums-and-errors:counting}
Write $|T|$ for the number of values of a finite type $T$. Then $|(A, B)| = |A| \cdot |B|$, and for an enum whose
variants carry the types $A_1, \dots, A_k$ (a variant with no field carrying the unit type \code{()}, with one
value), $|E| = |A_1| + \dots + |A_k|$. In particular $|\code{Option<A>}| = |A| + 1$ and
$|\code{Result<A, B>}| = |A| + |B|$.
\end{proposition}

\begin{proof}
A pair is determined by its first component, $|A|$ choices, and independently its second, $|B|$ choices: $|A||B|$
pairs. A value of the enum is a variant $i$ and a value of $A_i$; two values with different variants differ, so the
sets of values of the variants are disjoint, and their sizes add. \code{Option<A>} has the variants \code{Some(A)} and
\code{None}, carrying \code{()}: $|A| + 1$. \code{Result<A, B>} has \code{Ok(A)} and \code{Err(B)}.
\end{proof}

A C \code{struct} of a \code{bool} and a \code{union} of an \code{int} and a \code{double} can be in $2 \cdot (2^{32} +
2^{64})$ states, a value of the flag times a member stored, of which the program means only $2^{32} + 2^{64}$; the other half are the cases where the flag
and the stored member disagree, which nothing prevents. A Rust enum has exactly the values meant.

\section{Pattern matching}\label{sec:b2:rust-enums-and-errors:match}

\begin{definition}[Pattern, pattern matching]\label{def:b2:rust-enums-and-errors:matching}
A \emph{pattern}\index{pattern} describes a shape of value: a literal (\code{0}), a range (\code{1..=9}), a variable
that binds whatever it matches (\code{x}), the wildcard \code{_}, a variant with patterns for its fields
(\code{Shape::Rect { w, h }}), a tuple of patterns, or alternatives \code{p | q}. \emph{Pattern
matching}\index{pattern matching}, \code{match v { p1 => e1, p2 if c => e2, ... }}, tries the patterns in order and
evaluates the expression of the first one that matches, with its variables bound; an arm may add a \emph{guard},
\code{if c}, a condition on the bound variables.
\end{definition}

\begin{definition}[Exhaustive match]\label{def:b2:rust-enums-and-errors:exhaustive}
A \code{match} is \emph{exhaustive}\index{exhaustive match} if every value of the matched type matches at least one of
its \omterm{def:b2:rust-enums-and-errors:matching}{patterns}, guards not counted. Rust accepts only exhaustive matches.
\end{definition}

\omcode{code/bachelor-2/07-rust-enums-and-errors/listings/nonexhaustive.rs}{9}{14}{Corners of a shape, forgetting
circles}\label{lst:b2:rust-enums-and-errors:nonexhaustive}

\omoutput[2-6,13-14]{code/bachelor-2/07-rust-enums-and-errors/outputs/nonexhaustive.txt}{The missing variant,
named.}\label{out:b2:rust-enums-and-errors:nonexhaustive}

Exhaustiveness is what makes a \omterm{def:b2:rust-enums-and-errors:adt}{sum type} safe to extend: add a variant, and the compiler lists every \code{match} that
must now handle it, unless the \code{match} ends with a wildcard \code{_}, which silently takes the new variant too
(\cref{exo:b2:rust-enums-and-errors:2}). Write wildcards only where a default is right for every future variant.

\section{Options and results}\label{sec:b2:rust-enums-and-errors:errors}

\begin{definition}[Option type]\label{def:b2:rust-enums-and-errors:option}
The \emph{option type}\index{option type} \code{Option<T>} of the standard library is the enum with the variants
\code{Some(T)} and \code{None}. It is the type of a value that may be absent, and replaces the null pointer and the
reserved \emph{no value} code.
\end{definition}

\omcode{code/bachelor-2/07-rust-enums-and-errors/listings/option_demo.rs}{7}{32}{A search that may find
nothing}\label{lst:b2:rust-enums-and-errors:option}

\omoutput{code/bachelor-2/07-rust-enums-and-errors/outputs/option_demo.txt}{Absences handled four ways: \code{match},
\code{if let}, \code{unwrap_or} and \code{?}.}\label{out:b2:rust-enums-and-errors:option}

The C version returns \code{NULL} for \emph{not found}, and nothing obliges the caller to test it. The first lookup
succeeds; the second dereferences \code{NULL}.

\omoutput[2-9,15-15]{code/bachelor-2/07-rust-enums-and-errors/outputs/run_null_crash.txt}{The same lookups in C
(\code{null_crash.c}): \code{NULL->pages} reads address 8.}\label{out:b2:rust-enums-and-errors:null}

In Rust the book cannot be used before the \code{Option} is opened: \code{find(...).pages} does not compile, since an
\code{Option<&Book>} has no field \code{pages}. Failures that come with a reason use the second enum.

\begin{definition}[Result type, error propagation]\label{def:b2:rust-enums-and-errors:result}
The \emph{result type}\index{result type} \code{Result<T, E>} is the enum with the variants \code{Ok(T)}, a success
carrying its value, and \code{Err(E)}, a failure carrying an error of type \code{E}. \emph{Error
propagation}\index{error propagation} is the return of a failure to the caller: the operator \code{?} applied to a
\code{Result} gives the value inside \code{Ok}, or returns \code{Err(e)} from the enclosing function at once,
converting \code{e} with \code{From::from} into that function's error type.
\end{definition}

The stations' log, \code{data/town/weather.csv}, has a line per day and station. Each line is parsed into a
\code{Reading} or an error that says what is wrong.

\omcode{code/bachelor-2/07-rust-enums-and-errors/listings/weather.rs}{20}{57}{One line, one
\code{Result}}\label{lst:b2:rust-enums-and-errors:parse}

\code{read_log} applies \code{parse_line} to every line, keeps the readings, records each rejected line with its
number and error, and sets aside a second reading for the same day and station. Its own errors, a file that cannot be
read or has no header, form a third enum, \code{LogError}, and \code{From<std::io::Error> for LogError} lets \code{?}
convert the error of \code{read_to_string} (both are shown in the weekend problem). Run on the log, and on a
misspelt file name:

\omoutput{code/bachelor-2/07-rust-enums-and-errors/outputs/run_weather.txt}{Every line accounted
for.}\label{out:b2:rust-enums-and-errors:weather}

\begin{omfigure}
\begin{tikzpicture}[font=\scriptsize, f/.style={draw, rounded corners=2pt, fill=black!4, minimum width=36mm, minimum height=7mm, align=center}]
\node[f] (m) at (0,0) {\texttt{main}\\ matches \texttt{Err(LogError::Io(e))}};
\node[f] (r) at (5.4,0) {\texttt{read\_log}\\ \texttt{read\_to\_string(path)?}};
\node[f] (p) at (10.8,0) {\texttt{std::fs}\\ fails: no such file};
\draw[ptr] (m.north east) to[bend left=20] node[above, text=black!60] {calls} (r.north west);
\draw[ptr] (r.north east) to[bend left=20] node[above, text=black!60] {calls} (p.north west);
\draw[ptr, omThm, thick] (p.south west) to[bend left=20] node[below, align=center] {\texttt{Err(io::Error)}} (r.south east);
\draw[ptr, omThm, thick] (r.south west) to[bend left=20] node[below, align=center] {\texttt{?}: \texttt{From::from}, returns\\ \texttt{Err(LogError::Io(e))}} (m.south east);
\begin{scope}[yshift=-2.9cm]
\node[f] (r2) at (0,0) {\texttt{read\_log}: \texttt{match}\\ records \texttt{Err(e)},\\ reads the next line};
\node[f] (p2) at (5.4,0) {\texttt{parse\_line}\\ \texttt{number(\textquotedbl{}temp\_c\textquotedbl{}, f[2])?}};
\node[f] (n2) at (10.8,0) {\texttt{number}\\ \texttt{Err(Missing(\textquotedbl{}temp\_c\textquotedbl{}))}};
\draw[ptr, omThm, thick] (n2.west) -- node[above] {returned} (p2.east);
\draw[ptr, omThm, thick] (p2.west) -- node[above] {\texttt{?}} node[below, align=center] {returns it} (r2.east);
\end{scope}
\end{tikzpicture}
\par\smallskip{\small Two errors travelling back. Above, the missing file: \code{?} converts the \code{io::Error} and
returns it from \code{read_log}, and \code{main} reports it. Below, a missing temperature: \code{?} in
\code{parse_line} skips the remaining fields, and \code{read_log}, which matches instead of using \code{?}, records the
error and reads the next line.}
\end{omfigure}

\begin{definition}[Panic]\label{def:b2:rust-enums-and-errors:panic}
A \emph{panic}\index{panic} stops the current thread on a condition the program declares impossible: an index out of
bounds, an \code{unwrap} of \code{None} or of \code{Err}, an explicit \code{panic!}. The values of the thread are
dropped, a message gives the place and the cause, and if the thread is the main one the process exits with status
101.
\end{definition}

The division of labour is the point. A \code{Result} is for failures that the program expects and handles: a bad line
in a file, a missing file, a server that does not answer. A \omterm{def:b2:rust-enums-and-errors:panic}{panic} is for bugs: a state the code was written to make
impossible. \code{unwrap} says \emph{this cannot fail}; on input data, it is usually wrong
(\cref{exo:b2:rust-enums-and-errors:6}). C++ \omterm{def:b2:cpp-classes-raii:exception}{exceptions} (\cref{ch:b2:cpp-classes-raii}) travel up invisibly through
any call; a \code{Result} is part of the signature, and each \code{?} marks a place where the function may return.

\section{Recursive types}\label{sec:b2:rust-enums-and-errors:recursive}

\begin{definition}[Recursive type, box]\label{def:b2:rust-enums-and-errors:recursive}
A \emph{recursive type}\index{recursive type} has a variant or a field whose type contains the type itself. Its
values must be finite, and its size must be known: the recursive field is therefore a \emph{box}\index{box (Rust)},
\code{Box<T>}, a pointer to a value of type \code{T} on the heap that owns it and frees it when dropped.
\end{definition}

Without the \omterm{def:b2:rust-enums-and-errors:recursive}{box}, an \code{Expr} containing two \code{Expr}s would be larger than itself. With it, an \code{Add} holds
two pointers. An arithmetic expression in one variable is a tree:

\omcodes{code/bachelor-2/07-rust-enums-and-errors/listings/expr.rs}{5-10,14-38}{An expression type, evaluated and
simplified}\label{lst:b2:rust-enums-and-errors:expr}

\omoutput{code/bachelor-2/07-rust-enums-and-errors/outputs/expr.txt}{$(x \cdot (2 + (-1)) + 0) \cdot (3x + 4 \cdot 0)$
and its simplified form, equal at $x = 5$.}\label{out:b2:rust-enums-and-errors:expr}

The values are integers modulo $2^{64}$ (\code{wrapping_add}, \code{wrapping_mul}), a ring in which the rules
$0 + e = e$, $0 \cdot e = 0$ and $1 \cdot e = e$ hold exactly; with \code{f64} the second would fail when $e$ is
infinite.

\begin{proposition}[Structural recursion terminates]\label{prop:b2:rust-enums-and-errors:structural}
A function on the values of a \omterm{def:b2:rust-enums-and-errors:recursive}{recursive type}, defined by a \code{match} on the variants in which every recursive call
is made on a field of the matched value, and whose other operations terminate, terminates on every value.
\end{proposition}

\begin{proof}
By structural induction on the value. On a variant with no recursive field there is no recursive call, and the other
operations terminate. On a variant with recursive fields $f_1, \dots, f_m$, assume the function terminates on each
$f_i$ (the induction hypothesis, legitimate since each $f_i$ is a strictly smaller part of a finite value); the body
makes finitely many calls on them, each terminating, plus terminating operations. Every value of a \omterm{def:b2:rust-enums-and-errors:recursive}{recursive type} is
finite, being built by finitely many constructions, so the induction covers every value.
\end{proof}

\code{eval} and \code{simplify} recurse only on \code{a} and \code{b}, the fields of the matched \code{Add} or
\code{Mul}: they terminate. The simplifier is useful only if it is right.

\begin{theorem}[The simplifier preserves the value]\label{thm:b2:rust-enums-and-errors:simplify}
For every expression \code{e} and every \code{x}, \code{e.clone().simplify().eval(x) == e.eval(x)}.
\end{theorem}

\begin{proof}
By structural induction on \code{e}. Write $\llbracket e \rrbracket$ for \code{e.eval(x)} and $S(e)$ for the
simplified expression. For \code{Num} and \code{X}, $S(e) = e$. For \code{Add(a, b)}, the induction hypothesis gives
$\llbracket S(a) \rrbracket = \llbracket a \rrbracket$ and $\llbracket S(b) \rrbracket = \llbracket b \rrbracket$,
and $\llbracket e \rrbracket = \llbracket a \rrbracket + \llbracket b \rrbracket$. The first arm returns
\code{Num(m + n)} when $S(a) = \code{Num(m)}$ and $S(b) = \code{Num(n)}$, whose value is $m + n = \llbracket S(a)
\rrbracket + \llbracket S(b) \rrbracket$. The second returns $S(b)$ when $S(a) = \code{Num(0)}$, of value $0 +
\llbracket S(b) \rrbracket$, or symmetrically. The third returns \code{Add(S(a), S(b))}, of value $\llbracket S(a)
\rrbracket + \llbracket S(b) \rrbracket$. In every arm the value is $\llbracket a \rrbracket + \llbracket b
\rrbracket$. The \code{Mul} case is the same with the rules $0 \cdot v = v \cdot 0 = 0$ and $1 \cdot v = v \cdot 1 = v$
of the ring (\cref{exo:b2:rust-enums-and-errors:11}).
\end{proof}

The test crate \code{tests/enums} checks the theorem on 3000 random expressions, and checks that simplifying twice
changes nothing. A binary search tree is the same kind of type; \code{bst.rs} inserts the page counts of seven books.

\omcode{code/bachelor-2/07-rust-enums-and-errors/listings/bst.rs}{2}{21}{A tree that owns its
subtrees}\label{lst:b2:rust-enums-and-errors:bst}

\omoutput{code/bachelor-2/07-rust-enums-and-errors/outputs/bst.txt}{The in-order walk is
sorted.}\label{out:b2:rust-enums-and-errors:bst}

\begin{omfigure}
\begin{tikzpicture}[font=\scriptsize, n/.style={draw, rounded corners=2pt, fill=omDef!10, minimum width=9mm, minimum height=5mm},
    lf/.style={draw=black!40, fill=black!5, text=black!50, inner sep=1.5pt, font=\tiny},
    level distance=10mm, level 1/.style={sibling distance=46mm}, level 2/.style={sibling distance=24mm},
    level 3/.style={sibling distance=12mm}, edge from parent/.style={draw, ptr}]
\node[n] {189}
  child {node[n] {65}
    child {node[lf] {Leaf}}
    child {node[n] {133}
      child {node[n] {98}}
      child {node[lf] {Leaf}}}}
  child {node[n] {251}
    child {node[lf] {Leaf}}
    child {node[n] {412}
      child {node[n] {320}}
      child {node[lf] {Leaf}}}};
\end{tikzpicture}
\par\smallskip{\small The tree of \cref{out:b2:rust-enums-and-errors:bst}, built by inserting 189, 65, 251, 133, 412,
98 and 320. Each arrow is a \code{Box}: the node owns its children, and dropping the root frees the whole tree. The
leaves under 98 and 320 are not drawn.}
\end{omfigure}

\section{Layout}\label{sec:b2:rust-enums-and-errors:layout}

An enum is stored as a \emph{tag}, a small integer that says which variant, followed by room for the largest
variant's fields. The compiler may do better when a field type has values that can never occur.

\omcode{code/bachelor-2/07-rust-enums-and-errors/listings/sizes.rs}{11}{25}{The sizes of a few
types}\label{lst:b2:rust-enums-and-errors:sizes}

\omoutput{code/bachelor-2/07-rust-enums-and-errors/outputs/sizes.txt}{Bytes per value; the target has 8-byte
pointers.}\label{out:b2:rust-enums-and-errors:sizes}

\begin{omfigure}
\begin{tikzpicture}[font=\scriptsize, x=3.7mm]
\foreach \name/\y/\cells in {
  {Option<u32>}/0/{{4/tag/black!8},{4/u32/omDef!10}},
  {Shape}/-1/{{8/tag, padding/black!8},{8/f64/omDef!10},{8/f64/omDef!10},{8/f64/omDef!10}},
  {Option<\&u32>}/-2/{{8/address; None = 0/omThm!10}}} {
  \node[anchor=east, font=\scriptsize\ttfamily] at (-0.3,\y) {\name};
  \def\xx{0}
  \foreach \w/\lab/\col in \cells {
    \draw[fill=\col] (\xx,\y-0.3) rectangle ++(\w,0.6);
    \node[font=\tiny] at (\xx+\w/2,\y) {\lab};
    \pgfmathparse{\xx+\w}\xdef\xx{\pgfmathresult}
  }
}
\foreach \b in {0,8,16,24,32} \node[text=black!50, font=\tiny] at (\b,-2.6) {\b};
\node[text=black!50, font=\tiny, anchor=west] at (33,-2.6) {bytes};
\end{tikzpicture}
\par\smallskip{\small Three layouts from \cref{out:b2:rust-enums-and-errors:sizes}. \code{Option<u32>} needs a tag,
since all $2^{32}$ bit \omterm{def:b2:rust-enums-and-errors:matching}{patterns} of a \code{u32} are values; padded to the alignment of the \code{u32}, 8 bytes.
\code{Shape} holds a tag and its largest variant, the triangle's three sides. \code{Option<&u32>} needs no tag: a
reference is never null, so the address 0 is free to mean \code{None}.}
\end{omfigure}

% ledger: b2:option-npo
The last layout is a guarantee, not an accident: the standard library documents that \code{Option<T>} has the same
size as \code{T} when \code{T} is a reference, a \code{Box}, a function pointer or a non-zero integer, with
\code{None} represented by all zero bits. An optional reference therefore costs exactly what a C pointer that may be
\code{NULL} costs, and the check the C programmer may forget is required by the type. The same reasoning gives
\code{Option<bool>} one byte: a \code{bool} uses 2 of the 256 \omterm{def:b2:rust-enums-and-errors:matching}{patterns} of its byte, and \code{None} takes a third.

\begin{inpractice}[Error types in practice]
Give each module its own error enum, with one variant per failure the caller may want to distinguish, and data (a
line number, a file name) for the message. Implement \code{std::fmt::Display} for the message and \code{From} for the
errors of the lower layers, so that \code{?} converts them. Keep \code{unwrap} and \code{expect} for conditions that
are bugs if false, and say why in the \code{expect} message.
\end{inpractice}

\section{Exercises}\label{sec:b2:rust-enums-and-errors:exercises}

\begin{exercise}[$\star$]\label{exo:b2:rust-enums-and-errors:1}
Predict the five lines, then check them in the solution. What changes if the arm \code{0 => "zero"} is moved just
after \code{1..=9}? And if it is moved after the wildcard \code{_}?
\omcode{code/bachelor-2/07-rust-enums-and-errors/listings/exo_match.rs}{2}{22}{\omterm{def:b2:rust-enums-and-errors:matching}{Patterns}, guards and tuples}
\end{exercise}

\begin{exercise}[$\star$]\label{exo:b2:rust-enums-and-errors:2}
The library added the variant \code{Lost} to \code{LoanState}, and the program no longer compiles. Fix
\code{fine_cents} (a lost book costs 2000 cents). Then explain what would have happened if the last arm had been
\code{_ => 0}.
\omcode{code/bachelor-2/07-rust-enums-and-errors/listings/exo_new_variant.rs}{3}{16}{A match written before
\code{Lost}}
\end{exercise}

\begin{exercise}[$\star$]\label{exo:b2:rust-enums-and-errors:3}
Write \code{checked_div(a: i32, b: i32) -> Option<i32>}, which returns \code{None} when the quotient does not exist
or does not fit, and \code{mean(xs: &[f64]) -> Option<f64>}. In which two cases does \code{a / b} \omterm{def:b2:rust-enums-and-errors:panic}{panic}
in Rust?
\end{exercise}

\begin{exercise}[$\star$]\label{exo:b2:rust-enums-and-errors:4}
Using \cref{prop:b2:rust-enums-and-errors:counting}, count the values of: \code{(bool, bool, u8)};
\code{Option<Option<bool>>}; \code{Result<bool, Option<u8>>}; an enum with the variants \code{Red}, \code{Amber},
\code{Green} and \code{Flashing(bool)}; and \code{[bool; 2]}. Which of them could share the same set of values, up to
renaming?
\end{exercise}

\begin{exercise}[$\star\star$]\label{exo:b2:rust-enums-and-errors:5}
For each of the three calls, say where \code{total} returns and which lines are printed, then check against the
output in the solution.
\omcode{code/bachelor-2/07-rust-enums-and-errors/listings/exo_question.rs}{2}{22}{Two \code{?} in a row}
\end{exercise}

\begin{exercise}[$\star\star$]\label{exo:b2:rust-enums-and-errors:6}
The program stops at the third line. Explain the message, and rewrite the loop so that a bad line is reported with
its number and skipped, using \code{Option} and \code{Result} instead of \code{unwrap}.
\omcode{code/bachelor-2/07-rust-enums-and-errors/listings/exo_unwrap.rs}{8}{16}{Two \code{unwrap}s}
\omoutput[2-6]{code/bachelor-2/07-rust-enums-and-errors/outputs/exo_unwrap.txt}{The \omterm{def:b2:rust-enums-and-errors:panic}{panic}.}
\end{exercise}

\begin{exercise}[$\star\star$]\label{exo:b2:rust-enums-and-errors:7}
Write an enum \code{Json} for JSON-like values (null, booleans, numbers, strings, arrays, objects with keys in
order) and a method \code{to_text} that prints a value on one line, escaping \code{"}, the backslash and the newline
in strings. Test it on a book record whose title contains quotation marks.
\end{exercise}

\begin{exercise}[$\star\star$]\label{exo:b2:rust-enums-and-errors:8}
The mirror image of a tree swaps the left and right subtrees of every node. Define \code{mirror} on the \code{Tree}
of \cref{lst:b2:rust-enums-and-errors:bst} and prove by structural induction that \code{mirror(mirror(t)) == t} for
every tree. Prove also that the in-order walk of \code{mirror(t)} is the in-order walk of \code{t} reversed.
\end{exercise}

\begin{exercise}[$\star\star$]\label{exo:b2:rust-enums-and-errors:9}
Write \code{enum List<T> { Nil, Cons(T, Box<List<T>>) }} with \code{from_vec}, \code{len}, \code{rev} and
\code{map}. Write \code{len} and \code{rev} with loops: why does \code{rev} need no clone of the elements? Test that
\code{l.rev().rev() == l}.
\end{exercise}

\begin{exercise}[$\star\star\star$]\label{exo:b2:rust-enums-and-errors:10}
Extend the expression type to \code{f64} with \code{Sin} and \code{Cos}, and write \code{derive}, which returns the
derivative by the sum, product and chain rules, and \code{simplify}. Test \code{derive} on 500 random expressions
against the central difference $(f(x + h) - f(x - h)) / 2h$, with a tolerance you justify, and check that
simplification never makes an expression larger.
\end{exercise}

\begin{exercise}[$\star\star\star$]\label{exo:b2:rust-enums-and-errors:11}
(a) Prove that \code{eval} terminates on every expression with \cref{prop:b2:rust-enums-and-errors:structural}, and
bound the number of calls by the number of nodes. (b) Complete the proof of
\cref{thm:b2:rust-enums-and-errors:simplify}: the \code{Mul} case, arm by arm. (c) Which arm would break the theorem
if the values were \code{f64}? Give an expression.
\end{exercise}

\begin{exercise}[$\star\star\star$]\label{exo:b2:rust-enums-and-errors:12}
A configuration file has lines \code{key = value}, comments after a hash sign, and blank lines. Write
\code{parse(text) -> Result<BTreeMap<String, String>, ConfigError>}, where \code{ConfigError} carries the line number
and a \code{Kind}: no \code{=}, empty key, a key set twice, a value that starts like a number but is not one. Implement
\code{Display} so that \code{e.to_string()} gives messages such as \code{line 3: a set twice}.
\end{exercise}

\section{Problem: Reading the weather station logs}\label{sec:b2:rust-enums-and-errors:problem}

\begin{problem}[{Weekend problem --- Reading the weather station logs}]\label{pb:b2:rust-enums-and-errors:1}
The town's two stations, Hilltop and Riverside, log a line a day each for 2025. The log is known to hold defects: some
temperatures are missing, some are the stations' code $-99.0$ for \emph{no measurement}, and a line may be repeated.
This problem reads it with the types of this chapter and checks that no defect is accepted silently.

\needspace{4\baselineskip}\medskip\noindent\textbf{Part I --- Types for readings.}\par\nopagebreak
\begin{enumerate}[beginpenalty=10000]
\item Why is \code{Station} an enum rather than a \code{String}? What does deriving \code{PartialOrd} and \code{Ord} give it?
\item \code{Reading} is a product of five fields. Why is \code{temp_c} an \code{f64} and not an
\code{Option<f64>}?
\item List the variants of \code{LineError} and the line defect each one describes.
\item How many values has \code{Station}? And \code{Option<Station>}? How many bytes would you expect for each?
\item Why is \code{-99.0} detected in \code{number} rather than by the caller?
\item What would a C version of \code{parse_line} return, and how would the caller tell the five errors apart?
\end{enumerate}

\needspace{4\baselineskip}\medskip\noindent\textbf{Part II --- Parsing with \code{Result}.}\par\nopagebreak
\omcodes{code/bachelor-2/07-rust-enums-and-errors/listings/weather.rs}{58-68,76-93}{The file-level errors and
\code{read_log}}
\begin{enumerate}[resume, beginpenalty=10000]
\item In \code{parse_line}, which fields are still read when \code{temp_c} is missing?
\item \code{read_to_string(path)?} returns a \code{std::io::Error}, yet \code{read_log} returns \code{LogError}. Which
line makes this compile?
\item Why does \code{read_log} use \code{match} on \code{parse_line} instead of \code{?}?
\item The first \code{Ok} arm has a guard. Explain it, and why its order with the second \code{Ok} arm matters.
\item Which line of \cref{out:b2:rust-enums-and-errors:weather} comes from \code{LogError::Io}, and through how many
functions did the error travel?
\item Add a variant for a date that is not of the form \code{YYYY-MM-DD}. Which \code{match}es must change?
\end{enumerate}

\needspace{4\baselineskip}\medskip\noindent\textbf{Part III --- A tree of monthly aggregates.}\par\nopagebreak
\begin{enumerate}[resume, beginpenalty=10000]
\item Store, for each pair (station, month), the number of readings, their sum, minimum and maximum, in a
\code{BTreeMap<(Station, u32), Agg>}. Why does a tuple key with a derived order give the months of Hilltop first?
\item Use the \code{entry} API, \code{and_modify} and \code{or_insert}, to update the tree in one lookup per reading.
\item How many keys does the tree have, and how many readings in all?
\item Riverside has 31 lines in June 2025. How many readings does its June aggregate count, and why?
\item Which months lose a reading to a missing temperature or an error code?
\item Why must the mean be computed from the sum and the count rather than kept as a running mean?
\end{enumerate}

\needspace{4\baselineskip}\medskip\noindent\textbf{Part IV --- Every defect accounted for.}\par\nopagebreak
\begin{enumerate}[resume, beginpenalty=10000]
\item Check that the counts of \cref{out:b2:rust-enums-and-errors:weather} add up to the 731 lines of data.
\item Which test in \code{tests/enums} would fail if a missing temperature were read as 0.0?
\item How would the C function of the opening treat the two error codes? And the missing fields?
\item The duplicate is identical to the earlier line. Should it be rejected or merged? Justify.
\item A new defect type appears next year: a temperature written with a comma, \code{3,9}. What happens to such a
line with the current code, and is it accepted silently?
\item State the result of the problem: every planted defect found, with its count, and none accepted silently.
\end{enumerate}
\end{problem}
